% ==========================================
% PAGINA 3: INDICE DEI CONTENUTI
% ==========================================
\noindent{\small \textcolor{gray}{Git Happens}}
\tableofcontents

\section{Introduzione}

\subsection{Scopo del documento}
Il presente documento ha lo scopo primario di definire formalmente la pianificazione, l'organizzazione e il controllo delle attività necessarie al gruppo \textbf{Git Happens} per la realizzazione del progetto \textbf{Sonar}, commissionato dall'azienda proponente \textbf{Halue S.r.l. Società Benefit} nell'ambito del corso di Ingegneria del Software. \\
\\
In conformità con i processi delineati nello standard ISO/IEC 12207:1995 e recepiti nel documento \textit{Norme di Progetto}, la trattazione articola in modo strutturato gli aspetti gestionali e operativi del ciclo di vita del software:
\begin{itemize}
    \item l'identificazione, la valutazione e le strategie di mitigazione dei rischi di progetto;
    \item la definizione del modello di sviluppo adottato, la suddivisione dei ruoli e le regole di rotazione delle risorse all'interno del team;
    \item la pianificazione temporale delle attività, scandita per periodi di lavoro (sprint) e orientata al raggiungimento delle baseline di progetto (\textit{Requirements and Technology Baseline}\textsubscript{G} e \textit{Product Baseline}\textsubscript{G});
    \item la preventivazione delle ore e dei costi associati a ciascun ruolo e periodo, nonché la rendicontazione a consuntivo con relativa analisi degli scostamenti.
\end{itemize}
Considerata la natura iterativa e incrementale del processo di sviluppo, il \textit{Piano di Progetto} non è da intendersi come un artefatto statico, bensì come uno strumento dinamico di governo e monitoraggio, soggetto a revisione periodica e aggiornamento al termine di ciascuno sprint per riflettere l'effettivo stato di avanzamento dei lavori.

\subsection{Glossario}
Per prevenire incomprensioni legate all'uso di termini tecnici o specifici del dominio, il gruppo ha redatto il \textit{Glossario}. La presenza della lettera \textbf{G} posta a pedice di una parola indica che il termine corrisponde a una voce presente ed esaustivamente descritta all'interno di tale documento.


\subsection{Riferimenti}

\subsubsection{Riferimenti normativi}
\begin{itemize}
    \item \textit{Norme di Progetto}
    \item \textbf{Capitolato C8 - Sonar}: Documentazione vivente di uno store Shopify (Halue S.r.l. SB):\\
    \url{https://www.math.unipd.it/~tullio/IS-1/2026/Progetto/C8-Sonar.pdf}\\
    (Consultato: )
    \item \textbf{Regolamento del progetto didattico:} \\
    \url{https://www.math.unipd.it/~tullio/IS-1/2026/Dispense/PD1.pdf} (Consultato: );
    \item \textbf{Standard ISO/IEC 12207:1995:} \\
    \url{https://www.math.unipd.it/~tullio/IS-1/2009/Approfondimenti/ISO_12207-1995.pdf} (Consultato: ).
\end{itemize}

\subsubsection{Riferimenti informativi}
\begin{itemize}
    \item \textit{Glossario};
    \item \textbf{Slide del corso di Ingegneria del Software - Ciclo di vita del software:} \\
    \url{https://www.math.unipd.it/~tullio/IS-1/2026/Dispense/T02.pdf} (Consultato: );
    \item \textbf{Slide del corso di Ingegneria del Software - Gestione di progetto:} \\
    \url{https://www.math.unipd.it/~tullio/IS-1/2026/Dispense/T04.pdf} (Consultato: );
    \item \textbf{Documentazione Atlassian Jira:} \\
    \url{https://support.atlassian.com/jira-software-cloud/} (Consultato: );
    \item \textbf{Documentazione GitHub:} \\
    \url{https://docs.github.com/} (Consultato: );
    \item \textbf{Documentazione \LaTeX:} \\
    \url{https://www.latex-project.org/help/documentation/} (Consultato: ).
\end{itemize}

